\section{Tensor-factor maps and controlled partial traces}

\begin{definition}[Reindexing a complex matrix algebra]
    \label{def:matrix_reindex_star_equiv}
    \lean{Matrix.reindexStarAlgEquiv, Matrix.reindexStarAlgEquiv_apply,
        Matrix.symm_reindexStarAlgEquiv,
        Matrix.reindexStarAlgEquiv_trans_reindexStarAlgEquiv}
    \leanok
    For a bijection $e:A\simeq A'$, let
    $R_e(X)_{a',b'}=X_{e^{-1}(a'),e^{-1}(b')}$.
    For finite $A$ and $A'$, this is a complex star-algebra isomorphism,
    with inverse $R_{e^{-1}}$,
    and $R_f\circ R_e=R_{f\circ e}$.
\end{definition}

\begin{theorem}[Bipartite coefficients under product reindexing]
    \label{thm:bipartite_coefficients_reindex}
    \lean{Matrix.bipartiteSlice_reindex_prod, Matrix.operatorBlock_reindex_prod}
    \leanok
    Let $X$ be a complex matrix indexed by $A\times B$, and let
    $e:A\simeq A'$ and $f:B\simeq B'$ be bijections. Define matrices on the
    first factor by $S_{ij}(X)_{ab}=X_{(a,i),(b,j)}$, and matrices on the
    second factor by $T_{ab}(X)_{ij}=X_{(a,i),(b,j)}$. Then
    \begin{align}
        S_{ij}(R_{e\times f}(X))
         & =R_e(S_{f^{-1}(i),f^{-1}(j)}(X)), \notag \\
        T_{ab}(R_{e\times f}(X))
         & =R_f(T_{e^{-1}(a),e^{-1}(b)}(X)). \notag
    \end{align}
    Neither identity requires finite index sets. The second identity holds
    for matrices with entries in any type.
\end{theorem}

\begin{proof}\leanok
    Both identities follow by evaluating the two free indices: simultaneous
    reindexing replaces each coordinate by its inverse image.
\end{proof}

\begin{definition}[Map on the second tensor factor]
    \label{def:id_tensor_map}
    \lean{Matrix.idTensorMap}
    \leanok
    For a linear map $\Psi_B$ and a bipartite matrix $\rho_{AB}$, define
    \begin{align}
        (\id_A\otimes\Psi_B)(\rho_{AB})
         & =\Sigma_{B'A}\bigl((\Psi_B\otimes\id_A)
        (\Sigma_{AB}(\rho_{AB}))\bigr),
        \notag
    \end{align}
    where $\Sigma$ denotes the canonical exchange of the two tensor factors.
\end{definition}

\begin{definition}[Independent maps on both tensor factors]
    \label{def:tensor_map_both}
    \lean{Matrix.tensorMapBoth}
    \leanok
    \uses{def:id_tensor_map}
    For linear maps $\Phi_A$ and $\Psi_B$, set
    \begin{align}
        (\Phi_A\otimes\Psi_B)(\rho)
         & =(\id_{A'}\otimes\Psi_B)
        ((\Phi_A\otimes\id_B)(\rho)).
        \notag
    \end{align}
\end{definition}

\begin{lemma}[Positivity under a map on the first tensor factor]
    \label{lem:tensor_map_id_pos}
    \lean{IsKrausCPTP.tensorMapId_posSemidef}
    \leanok
    If $\Phi_A$ is trace-preserving completely positive and
    $\rho_{AB}\geq0$, then
    $(\Phi_A\otimes\id_B)(\rho)\geq0$.
\end{lemma}

\begin{lemma}[Positivity under a map on the second tensor factor]
    \label{lem:id_tensor_map_pos}
    \lean{IsKrausCPTP.idTensorMap_posSemidef}
    \leanok
    \uses{def:id_tensor_map, lem:tensor_map_id_pos}
    If $\Psi_B$ is trace-preserving completely positive and
    $\rho_{AB}\geq0$, then
    $(\id_A\otimes\Psi_B)(\rho)\geq0$.
\end{lemma}

\begin{lemma}[Positivity under independent maps]
    \label{lem:tensor_map_both_pos}
    \lean{IsKrausCPTP.tensorMapBoth_posSemidef}
    \leanok
    \uses{def:tensor_map_both, lem:tensor_map_id_pos, lem:id_tensor_map_pos}
    If $\Phi_A$ and $\Psi_B$ are trace-preserving completely positive and
    $\rho_{AB}\geq0$, then $(\Phi_A\otimes\Psi_B)(\rho)\geq0$.
\end{lemma}
\begin{definition}[Controlled dependent partial trace]
    \label{def:controlled_dependent_partial_trace}
    \lean{Matrix.controlledPartialTraceRightLM}
    \leanok

    Let
    \begin{align}
        H & =\bigoplus_{i\in I}(A_i\otimes B_i), \notag \\
        K & =\bigoplus_{i\in I}A_i. \notag
    \end{align}
    The controlled dependent partial trace discards the off-diagonal blocks
    between distinct $i$ and applies $\tr_{B_i}$ to the $i$th diagonal block.
\end{definition}

\begin{theorem}[Diagonal blocks of the controlled dependent partial trace]
    \label{thm:controlled_dependent_partial_trace_same_block}
    \lean{Matrix.controlledPartialTraceRightLM_sameBlock_apply}
    \leanok
    \uses{def:controlled_dependent_partial_trace}
    If $X_{ii}$ denotes the $i$th diagonal block of $X$, then
    $[\mathcal C_{\tr}(X)]_{ii}=\tr_{B_i}(X_{ii})$.
\end{theorem}
\begin{proof}
    \leanok

    If $(K_{i,j})_j$ is the Kraus family chosen for $\tr_{B_i}$, then the
    diagonal-block formula gives
    \begin{align}
        [\mathcal C_{\tr}(X)]_{ii}
         & =\sum_j K_{i,j}X_{ii}K_{i,j}^\dagger
        =\tr_{B_i}(X_{ii}). \notag
    \end{align}
\end{proof}

\begin{theorem}[The controlled dependent partial trace is trace-preserving
        completely positive]
    \label{thm:controlled_dependent_partial_trace_cptp}
    \lean{Matrix.controlledPartialTraceRightLM_isKrausCPTP}
    \leanok
    \uses{def:controlled_dependent_partial_trace}
    The controlled dependent partial trace is trace-preserving and completely
    positive.
\end{theorem}
\begin{proof}
    \leanok

    Choose Kraus operators $(V_{i,a})_a$ for $\tr_{B_i}$ and extend each
    $V_{i,a}$ by zero outside the $i$th summand. Their controlled family
    satisfies
    \begin{align}
        \sum_{i,a}\widetilde V_{i,a}^{\dagger}\widetilde V_{i,a}
         & =\bigoplus_i\left(\sum_aV_{i,a}^{\dagger}V_{i,a}\right)
        =\bigoplus_i \Id_{A_i\otimes B_i}
        =\Id_H. \notag
    \end{align}
    Its Kraus form gives complete positivity, and the displayed resolution of
    the identity gives trace preservation.
\end{proof}
